% -*- coding: utf-8 -*-
% !TEX program = xelatex
\documentclass[plain,noanswer,medmath]{../../template/jnuexam}
\usepackage{../../template/kysx}

\begin{document}


\examtitle{1990}{3}


\loadproblems{1}{1990P1}% 载入试卷一
\loadproblems{2}{1990P2}% 载入试卷二


\makepart{填空题}{本题满分15分，每小题3分}


\begin{problem}
曲线$\begin{cases}
  x= \cos^3 t \\
  y= \sin^3 t
\end{cases}$上对应于点$t=\frac{\pi}{6}$点处的法线方程是\fillin{}.
\end{problem}


\begin{problem}
设$y=\e^{\tan\frac{1}{x}}\sin\frac{1}{x}$，则$y'=$\fillin{}.
\end{problem}


\begin{problem}
$\int_0^1x\sqrt{1-x}\dx=$\fillin{}.
\end{problem}


\begin{problem}
下列两个积分的大小关系是：$\int_{-2}^{-1}\e^{-x^3}\dx$\fillin{}$\int_{-2}^{-1}\e^{x^3}\dx$．
\end{problem}


\useproblem{1}{1}{3}%【同试卷一第一(3)题】


\makepart{选择题}{本题满分15分，每小题3分}


\begin{problem}
已知$\lim_{x\to \infty}\,\left(\frac{x^2}{x+1}-ax-b \right)=0$，其中$a,b$是常数，则\pickout{}
\begin{abcd}
\item $a=1,b=1$．
\item $a=-1,b=1$．
\item $a=1,b=-1$．
\item $a=-1,b=-1$．
\end{abcd}
\end{problem}


\begin{problem}
设函数$f(x)$在$(-\infty +\infty)$上连续，则$\d\left[\int f(x)\dx \right]$等于\pickout{}
\begin{abcd}
\item $f(x)$．
\item $f(x)\dx$．
\item $f(x)+C$．
\item $f'(x)\dx$．
\end{abcd}
\end{problem}


\useproblem{1}{2}{2}%【同试卷一第二(2)题】


\useproblem{1}{2}{1}%【同试卷一第二(1)题】


\begin{problem}
设$F(x)=\begin{cases}
  \frac{f(x)}{x}, & x \ne 0 \\
            f(0), & x = 0
\end{cases}$，其中$f(x)$在$x=0$处可导，$f'(0)\ne 0,f(0)=0$，
则$x=0$是$F(x)$的\pickout{}
\begin{abcd}
\item 连续点．
\item 第一类间断点．
\item 第二类间断点．
\item 连续点或间断点不能由此确定．
\end{abcd}
\end{problem}


\makepart{计算题}{本题满分25分，每小题5分}


\begin{problem}
已知$\lim_{x\to\infty}\left(\frac{x+a}{x-a}\right)^x=9$，求常数$a$．
\end{problem}


\begin{problem}
求由方程$2y-x=(x-y)\ln(x-y)$所确定的函数$y=y(x)$的微分$\dy$．
\end{problem}


\begin{problem}
求曲线$y=\frac{1}{1+x^2}$（$x>0$）的拐点．
\end{problem}


\begin{problem}
计算$\int \frac{\ln x}{(1-x)^2} \dx.$
\end{problem}


\useproblem{2}{4}{2}%【同试卷二第四(2)题】


\makepart{}{本题满分9分}

\begin{problem}
在椭圆$\frac{x^2}{a^2}+\frac{y^2}{b^2}=1$的第一象限部分上求一点$P$，使该点处的切线、椭圆及两坐标轴所围图形面积为最小（其中$a>0,b>0$）．
\end{problem}


\makepart{}{本题满分9分}

\begin{problem}
证明：当$x>0$时，有不等式$\arctan x+\frac{1}{x}>\frac{\pi}{2}$．
\end{problem}


\makepart{}{本题满分9分}

\begin{problem}
设$f(x)=\int_1^x \frac{\ln t}{1+t} \dt$，其中$x>0$，求$f(x)+f\left(\frac{1}{x} \right)$．
\end{problem}


\makepart{}{本题满分9分}%【同试卷二第四(3)题】

\useproblem{2}{4}{3}


\makepart{}{本题满分8分}

\begin{problem}
求微分方程$y''+4y'+4y= \e^{ax}$之通解，其中$a$为实数．
\end{problem}


\end{document}
